%% Plain LaTeX article. No IEEEtran anywhere.
%% Body text is uniform 10pt. Title and section headings are enlarged.
%% Abstract wording unchanged from the author.

\documentclass[10pt,a4paper]{article}

\usepackage[utf8]{inputenc}
\usepackage[T1]{fontenc}
\usepackage{lmodern}
\usepackage{microtype}

% ---------------------------------------------------------------------
% MARGINS
% One value for all four sides:
%     \usepackage[margin=2cm]{geometry}
% Separate values, which is what is active below. Later keys win, so
% "margin" sets a baseline and top/bottom override it.
% Increase top/bottom to push the text block inward, decrease to gain room.
% ---------------------------------------------------------------------
\usepackage[margin=2cm, top=2cm, bottom=2cm]{geometry}

\pagestyle{empty}

\hyphenation{pho-to-vol-ta-ic ir-ra-di-ance micro-con-trol-ler}

\begin{document}

% ---------- Title block ----------
\begin{center}
{\LARGE\bfseries Maximum Power Point Tracking for Photovoltaic\\[0.3em]
Arrays under Partial Shading Conditions\par}

\vspace{1em}
Soumyadip Das \quad B Kishan \quad Aditi Arora
\end{center}

\vspace{1.5em}

% ---------- Abstract. Plain paragraphs, no abstract environment. ----------
\begin{center}
{\large\bfseries Abstract}
\end{center}

\noindent
Reliable and efficient photovoltaic (PV) power generation is essential for
maximising the utilisation of renewable energy. However, conventional Maximum
Power Point Tracking (MPPT) techniques face significant challenges under
partial shading conditions, where non-uniform irradiance across PV modules
can result in multiple peaks in the power--voltage (P--V) characteristic.
This makes it difficult for conventional tracking techniques to consistently
identify and track the Global Maximum Power Point (GMPP), potentially leading
to significant power losses.

This project aims to investigate and develop an effective MPPT strategy for
photovoltaic arrays operating under partial shading conditions. A PV array
integrated with a DC--DC converter will be modelled and simulated to study its
behaviour under both uniform irradiance and different partial shading patterns.
Various MPPT approaches will be investigated and comparatively evaluated, with
conventional techniques such as Perturb and Observe (P\&O) and Incremental
Conductance (INC) studied alongside global optimisation and intelligent
techniques, including Genetic Algorithms (GA), Artificial Neural Networks
(ANNs), and other suitable approaches. The methods will be evaluated based on
parameters such as global maximum power point tracking accuracy, convergence
speed, dynamic response, steady-state oscillations, computational complexity,
and overall energy extraction. Based on the comparative analysis, a suitable
approach will be identified and its effectiveness under different operating
conditions will be assessed. The project ultimately aims to improve the
reliability and efficiency of PV power extraction under partial shading
conditions while establishing an effective balance between tracking performance
and computational requirements.

\vspace{1em}
\noindent
\textbf{Keywords:} artificial neural networks, DC--DC converters, genetic
algorithms, global maximum power point, maximum power point tracking, partial
shading, photovoltaic systems.

\vspace{1.5em}

% ---------- Objectives ----------
\begin{center}
{\large\bfseries Objectives}
\end{center}

\begin{enumerate}
\setlength{\itemsep}{6pt}
\setlength{\parskip}{0pt}

\item \textbf{Literature Study and Qualitative Comparison.}
To survey existing maximum power point tracking techniques for partially shaded
arrays and classify them by solution family, covering conventional
hill-climbing, scanning, model-based, metaheuristic, soft-computing, and hybrid
schemes. Each family will be compared on tracking mechanism, sensing
requirement, controller class, and determinism, with reported results separated
into simulation and hardware-validated evidence, and the survey extended to
commercial inverter practice and the applicable MPPT efficiency standards.

\item \textbf{Development and Simulation of a Robust GMPPT Algorithm.}
To formulate and develop a Global Maximum Power Point Tracking (GMPPT)
algorithm capable of reliably navigating the multi-peak P--V curves caused by
partial shading. The algorithm will be implemented on a simulated PV array
coupled to a DC--DC converter and exercised across uniform irradiance and a
representative set of shading patterns, with simulation studies used to
validate its convergence speed, transient response, and steady-state stability.

\item \textbf{Real-Time Hardware Validation.}
To implement the proposed algorithm on a rapid control prototyping platform
comprising a host PC running MATLAB/Simulink and a dSPACE MicroLabBox
real-time controller, driving a DC--DC boost converter under reproducible
shading. Measured tracking efficiency, convergence time, and the number of
operating points applied to the array will be compared against simulation, and
the execution time and memory footprint recorded to establish portability to a
low-cost microcontroller.

\item \textbf{Comparative Performance Analysis.}
To execute an empirical comparison of the proposed GMPPT strategy against
conventional derivative-based methods such as Perturb and Observe, and against
standard metaheuristic optimisation algorithms including Particle Swarm
Optimisation (PSO), Genetic Algorithms (GA), Ant Colony Optimisation (ACO), and
the Grey Wolf Optimizer (GWO). The comparison will benchmark dynamic tracking
efficiency, convergence time on the true Global Maximum Power Point (GMPP), and
steady-state power oscillation under complex shading profiles.

\end{enumerate}

\end{document}
